\begin{abstract}
Processing contexts far beyond a backbone's native window remains costly:
full attention grows quadratically and, past the extrapolation limit,
simply collapses. We introduce \textbf{QCMem} (Query-Conditioned Memory),
which recasts long-context processing from \emph{token-partial} KV reuse
(recompute the full-depth KV of a subset of tokens) to
\emph{layer-partial} readout: split the transformer at depth $j$, cache the
single depth-$j$ residual-stream hidden of each chunk, then at read time
retrieve the top-$k$ relevant chunks and recompute only the upper layers
$[j{:}L]$ over a fixed-size pack. The read cost and memory are independent of
context length (constant $\approx$6.7k tokens / $\approx$18\,GB for an 8B
backbone). Across five long-context benchmarks (RULER, BABILong, LongEval,
LongBench, LoCoMo) with three mechanism-level baselines under an identical
harness, QCMem matches full-context accuracy inside the window and is
\emph{the only method that still works past the backbone's extrapolation
limit}: at 128k, full-context / full-KV recompute collapse to 0 while QCMem
stays at RULER~100 and LongEval~0.98, with a $7.8\times$ prefill speedup and
constant memory. We attribute the effectiveness to an
understand-then-generate division of labor---semantic content saturates in
mid layers---and give both correlational probes and a causal end-to-end test.
\end{abstract}
